\chapter*{Introduction}
\addcontentsline{toc}{chapter}{Introduction}
\markboth{Introduction}{Introduction}
\label{chap:introduction}

\noindent In recent years, the integration of artificial intelligence and
machine learning has profoundly reshaped the way organisations manage
their human-capital decisions. With the rise of large language models,
multimodal vision and speech encoders and agentic frameworks, machines now
understand documents, conduct conversations and analyse behaviour with
unprecedented fluency. Yet, despite this progress, the application of
these techniques to recruitment, one of the most consequential and most
regulated decisions an organisation makes, remains fragmented, shallow
and largely opaque.

\vspace{0.3cm}

\noindent Hiring teams currently rely on a patchwork of disconnected
products: applicant-tracking systems perform keyword filtering, sourcing
platforms expose closed proprietary scores, and asynchronous video tools
record candidates without ever adapting to them. At the same time, the
European AI Act and the NIST AI Risk Management Framework now classify
employment-related AI as \emph{high-risk} and require a documented
decision trail and meaningful human oversight at every stage.

\vspace{0.3cm}

\noindent The goal of this project, titled \textbf{HireFlow}, is to design
and deliver an end-to-end agentic platform that supports a recruiter
throughout the full hiring journey, from a vacancy opening to the final
hiring decision. Built around five collaborating agents responsible for
culture-aware job-description authoring, candidate--job matching,
conversational interviewing, multimodal behavioural analysis and
recruiter-adaptive scoring, the system leverages transformer models,
graph neural networks and computer-vision pipelines while preserving
meaningful human oversight at every step.

\vspace{0.3cm}

\noindent This report is structured as follows. The first chapter
situates the project, reviews existing platforms and presents the
proposed multi-agent solution. The remaining chapters follow the Scrum
cadence of the internship --- one chapter per sprint --- and the sprints
that build the artificial-intelligence agents are themselves organised
around the six phases of CRISP-DM. The final sections discuss the
results, limitations and perspectives for future improvement.

